\documentclass[../main.tex]{subfiles}
\begin{document}
%==========================================TIÊU ĐỀ TẬP========================================
\begin{table}[h]
  \begin{adjustwidth}{-1cm}{}
    \begin{tabular}{>{\centering\arraybackslash}p{6cm}|>{\raggedright\arraybackslash}p{12.5cm}}
      \multirow{5}{6cm}
      % Cột bên trái: ÉPISODE và số tập
      {\\[-40pt]
        \flaregothic\fontsize{20pt}{20pt}\selectfont \centering
        {\contour{errorthree}{\textcolor{black}{ÉPISODE}}}\color{black}\\
        \burbank\fontsize{72pt}{72pt}\selectfont
        \includegraphics[height=3cm]{../../../../Image/Book?/number/num9.png}
        \\[18pt]
        % \flaregothic\fontsize{14pt}{14pt}\selectfont
        % \color{epattr}BIENVENUE, \uline{\color{Banpire} Mel} !
        % \flaregothic\fontsize{18pt}{18pt}\selectfont
        % \color{epattr}ÉPISODE PERDU
        \color{black}
      }
       &
      % Dòng thứ nhất: tên tập tiếng Nhật
      {\fotskip\fontsize{18pt}{18pt}\selectfont \ruby{青}{あお}\ruby{上}{がみ}\ruby{駅}{えき}の\ruby{亡}{ぼう}\ruby{霊}{れい}
      } \\[6pt]
      % Dòng thứ hai: tên tập tiếng Anh
       & {\cafeteria\fontsize{18pt}{18pt}\selectfont The Ghost of Aogami Station}                            \\[4pt]
      % Dòng thứ ba: tên tập tiếng Việt
       & {\vnmsans\fontsize{18pt}{18pt}\selectfont Bóng ma trong nhà ga Aogami}                              \\[8pt]
      % Dòng thứ tư: Nhân vật xuất hiện
       & {\caviar{\fontsize{14pt}{14pt}\selectfont Track used: The Crystal - Łukasz Michalski}}
      \\[6pt]
      % Dòng cuối cùng: Thời gian diễn ra của tập (thường sẽ là ngày phát hành)
       & {\continuum\fontsize{16pt}{16pt}\selectfont Ngày phát hành: 15/04/2022}                             \\[8pt]
    \end{tabular}
  \end{adjustwidth}
\end{table}
%===========================================NỘI DUNG==========================================
\color{black}

\begin{center}
  \uline{Tháng Hai, năm 2018.}
  \pov{Haragen Kyuen}
\end{center}

*LẠCH CẠCH* *LẠCH CẠCH*

*LẠCH CẠCH* *LẠCH CẠCH*

Tiếng bánh xe tàu nghiến trên đường ray vọng khắp toa, tạo nên âm thanh đều đặn như nhịp đập của trái tim sắt thép.

``\dots''

``\dots''

Hòa cùng với đó là tiếng loa phát thanh rè rè:

``Điểm dừng tiếp theo: ga Aogami. Xin quý khách chuẩn bị kiểm tra kỹ tư trang hành lý trước khi xuống tàu.''

*BENG BENG BENG BENG\dots*

Đâu đó xa xa, tôi nghe thấy tiếng đèn tín hiệu nhấp nháy, báo hiệu tàu đang đi qua.

Một thứ cảm giác quen thuộc, như thể những lúc buổi sáng tôi tới trường, hay buổi chiều giờ tan tầm.

``\dots sama\dots\ Onii-sama!''

Giọng một cô gái vang vọng lên ở đâu đó. Một giọng nói rất quen thuộc, mà tôi đã nghe thấy ở đâu rồi\dots

``\dots Ư\dots''

``Dậy thôi nào, Onii-sama, ta sắp tới nơi rồi!''

Tiếng rền của tàu kèm với tiếng thúc giục của cô ấy kéo tôi ra khỏi giấc ngủ chập chờn. Tôi đã ngủ từ khi nào? Và trong bao lâu?

``Ái da\dots''

Đầu tôi nặng như thể vừa nhét đầy một mớ hỗn loạn không thể sắp xếp nổi. Một cơn mơ kỳ lạ vẫn còn hằn in trong trí óc: ngôi trường sụp đổ, hành lang lạnh lẽo phủ bụi, hoa bỉ ngạn đỏ nở rực như báo hiệu cái chết. Tôi thấy những khuôn mặt quen thuộc, bạn học kêu gào rồi im bặt, rồi đến Shino với đôi mắt xanh sáng lạnh, gọi chúng tôi là ``Hikawa-san,'' và Sora - bản sao của cô ấy, với đôi mắt nâu đầy xa lạ, bên trong con người ấy là niềm khao khát kết bạn tới mãnh liệt.

Thậm chí còn có Reiko-san, với mái tóc đen và đôi mắt hai màu, đồng hành cùng chúng tôi trong quá khứ xa xăm. Mọi thứ như chồng chéo lên nhau, vừa thực vừa hư, rợn ngợp đến mức tôi rùng mình. Nếu chỉ là mơ\dots\ sao nó lại quá chân thật đến vậy?

``\dots''

``\dots''

\dots Có lẽ tôi thực sự phải ngồi dậy thôi.

Tôi từ từ mở mắt. Ánh sáng mờ mịt từ cửa sổ tàu hắt lên khuôn mặt người con gái ngồi bên cạnh. Mái tóc đen dài tới ngang vai được buộc hai bên, đôi mắt dị sắc xanh-đỏ quen thuộc\dots\ giống hệt nhân vật trong mơ. Chỉ khác là, hôm nay cô mặc đồng phục hiện đại với váy ca-rô xanh thẫm, cài một chiếc nơ cùng màu trước ngực, khoác áo len nâu bên ngoài. Trên tay trái của cô ấy có đeo một chiếc vòng màu đỏ với ba chiếc chuông vàng nhỏ đính chụm.

``Reiko-san\dots?'' tôi khẽ gọi, giọng khàn đục.

Người vừa gọi quay lại, nở nụ cười dịu dàng. ``Onii-sama, anh tỉnh rồi à? Em lo lắm, anh ngủ say đến mức tàu rung cũng chẳng hay à\~{}''

\dots Hả? ``Onii-sama?''

Tôi chớp mắt, khẽ giật mình. Tại sao Reiko-san lại gọi tôi như anh trai? Ở trường Aogami, cô ấy chỉ là bạn, một người mới quen mang chung họ Haragen. Nhưng giờ\dots

\dots Không lẽ nào, ký ức của tôi\dots\ lại đang rối loạn rồi sao?

``Tôi\dots\ vừa nghe nhầm chăng? Reiko-san, cậu vừa gọi tôi là Onii-sama sao?'' tôi lắp bắp hỏi, nhìn cô đầy ngờ vực.

Cô ấy nghiêng đầu, đôi mắt hai màu lấp lánh. ``Trời ạ\dots\ anh vẫn còn ngái ngủ sao? Onii-sama là Onii-sama, không phải anh thì ai chứ?''

Tôi cứng họng, chẳng biết phải đáp lại thế nào. Tim tôi đập càng gấp, còn đầu óc thì xoay vòng trong mớ hỗn loạn mới chồng thêm vào giấc mơ vừa rồi.

Từ khi nào cô ấy lại trở thành em gái tôi vậy? Tôi nhớ rằng chỉ có Saikyou là em gái tôi thôi chứ?

``Mà nhắc tới em gái\dots\ Saikyou đâu rồi?'' tôi vô thức cất lời.

Chưa để Reiko-san - mà có lẽ tôi nên bỏ hậu tố ``-san'' đi, vì tôi cảm tưởng rằng cô ấy sẽ không thích bị gọi vậy - nói thêm điều gì, cửa toa trượt mở, Saikyou bước vào với cốc nước trên tay. Em ấy nhìn tôi, nở một nụ cười hiền xen  lẫn sự nhẹ nhõm. ``Ái chà, Nii-chan, anh dậy rồi sao?''

``Onee-sama!'' Reiko reo lên, nhào tới ôm chầm lấy cô em.

Saikyou bật cười, đưa tay xoa đầu em: ``Reiko ngoan lắm, trông coi Nii-chan ngủ như thế này là tốt rồi\~{}'' rồi ngoảnh về tôi: ``Nii-chan, anh tỉnh rồi à? Trông anh như người mất hồn vậy.''

Bộ đồng phục của em ấy giống hệt Reiko, nhưng cô mặc thêm một áo blazer đen bên ngoài thay vì áo len.

Tôi trố mắt nhìn hai chị em ôm nhau, não tôi vẫn chưa kịp thích ứng nổi. ``Khoan\dots\ Saikyou, chúng ta là sinh đôi, sao lại\dots\ sao lại có thêm một đứa em gái nữa chứ!?''

Saikyou thoáng lúng túng, nhưng nhanh chóng lấy lại bình tĩnh. Em khẽ gật, như thể đã đoán được gì đó: ``Lúc em tỉnh lại cũng bất ngờ y hệt anh, em cũng không thể nghĩ là em sẽ có một đứa em gái. Nhưng khi Reiko gọi em là `Onee-sama'\dots chẳng hiểu sao em lại thấy quen thuộc, gần gũi. Cứ như vốn dĩ em đã là chị của con bé vậy.''

Reiko thì xoa trán, vẻ mặt hơi nhăn nhó, như đang mô tả lại cảm giác của mình khi tỉnh lại. ``Em cũng thế. Khi mở mắt, em cũng nhức đầu khủng khiếp. Nhưng em chỉ biết\dots\ em có một người anh và một người chị. Anh là Kyuen-onii-sama, chị là Saikyou-onee-sama.''

Tôi ngồi sững với câu trả lời có phần thoải mái của cô ấy. Vậy ra, không chỉ riêng tôi, mà cả hai người họ cũng bị cuốn vào trò xáo trộn ký ức quái gở này. Đó có phải là những gì xảy ra sau khi tôi bước vào cánh cổng đền đó không?

Nhìn hai người họ, tôi cảm thấy có gì đó rất đỗi thân thuộc. Cách Reiko nép vào Saikyou, cách em ấy xoa đầu Reiko\dots\ tất cả đều tự nhiên đến lạ. Cảm tưởng như thể họ vốn dĩ đã chị em từ rất lâu rồi vậy.

Cả cái cách mà Reiko thoải mái nói chuyện với tôi, và cách gọi hai chúng tôi là Onii-sama và Onee-sama\dots\ hình như tôi đã từng nghe thấy ở đâu đó\dots\ mà tôi chẳng tài nào nhớ nổi.

Chợt, một ý nghĩ lóe lên trong đầu tôi. Có điều gì đó rất quan trọng mà tôi cần phải biết ngay.

``Saikyou\dots\ giờ là năm nào?'' tôi hỏi, giọng run run.

Em tôi giơ điện thoại, màn hình sáng loáng. ``2018. Chính xác là năm tháng trước khi chúng ta nhập học.''

Năm tháng trước\dots\ nghĩa là đã quay về thời hiện đại. Nhưng lại là tháng Hai. Khoảng thời gian mà chúng tôi chuyển về thị trấn Aogami để bắt đầu sinh sống tại đây.

\dots Cơ mà, từ từ.

Nếu đây là thời điểm năm tháng trước khi chúng tôi nhập học, thì tại sao chúng tôi lại đang diện bộ đồng phục của ngôi trường đó? Tôi không hề nhớ rằng chúng tôi đã có bộ đồng phục đó cho tới tận ba ngày trước khi nhập học đâu?

Chưa kể rằng bộ đồng phục tôi đang mặc nữa - bộ áo ba lớp gồm áo sơ mi bên trong, áo len nâu nằm giữa và áo khoác blazer màu den bên ngoài cùng, ở giữa ngực là thắt một chiếc cà vạt xanh thay vì là cái nơ - dù rằng cũng ổn áp, nhưng đó không phải là điều tôi đang nghĩ tới lúc này!

``Này, Saikyou, em có thấy lạ không?'' tôi hỏi, giọng đầy nghi hoặc. ``Sao chúng ta lại mặc đồng phục của trường vậy? Nếu theo dòng thời gian thì tận năm tháng nữa mới nhập học mà?''

Em tôi nhíu mày, liếc xuống bộ cánh đang mặc, rồi khẽ giật mình như thể vừa nhận ra điều gì đó không ổn: ``À, phải rồi\dots\ Em cũng thấy kỳ lạ. Bình thường phải tới gần ngày nhập học mới có đồng phục chứ?''

Nhưng, Reiko thì vẫn đang mải mê ngắm cảnh qua cửa sổ, dường như không hề để tâm đến chi tiết kỳ lạ này. Cô chỉ khẽ đu đưa người theo nhịp tàu chạy, như thể việc mặc đồng phục là điều hiển nhiên nhất trên đời.

Tôi thở dài, lưng ngả xuống ghế tàu, để mặc cho tiếng lạch cạch đều đặn của tàu đưa suy nghĩ trôi dạt. Đầu óc mụ mị, tôi vẫn chẳng biết rốt cuộc đâu là giấc mơ, đâu là thực tại, mọi thứ trong đầu tôi lúc nào chẳng khác nào một mớ dây mơ rễ má.

Nhưng, có một điều rõ ràng: giờ đây tôi không chỉ có một cô em gái. Tôi có đến hai người, Saikyou và Reiko. Dù khó hiểu đến mức nào, tôi cũng chẳng còn lựa chọn nào khác ngoài chấp nhận thực tại này, ít nhất là cho tới khi tìm được lời giải.

\dots

\dots Khoảng thời gian này\dots\ hình như tôi đã từng trải qua rồi. Tháng Hai, chuyến tàu vào Aogami, cùng Saikyou, trong lúc bố mẹ tôi hồi đó đã vào trước để sắp xếp đồ đạc. Chúng tôi chuyển trường\dots\ lần thứ mấy rồi nhỉ? Tôi không tài nào nhớ nổi, và tôi cũng chẳng buồn để tâm. Cứ như có ai đó liên tục nhấn nút ``làm lại,'' ép hai anh em tôi lặp đi lặp lại một kịch bản mơ hồ.

``Nii-chan\dots'' khẽ thì thầm, như đọc được ý nghĩ trong tôi. ``Cảm giác này\dots\ đúng là quen thuộc. Nhưng em lại thấy có chút gì đó khác. Như thể cùng một bức tranh, nhưng màu sắc bị pha loãng.''

Tôi nhìn em gái mình, rồi lại liếc sang Reiko. Đối với Saikyou, đây là một sự từng trải, nhưng với cô ấy, tất cả đều mới mẻ, chí ít là khi tới Aogami với một diện mạo hoàn toàn khác xa so với thời kỳ hậu chiến.

Reiko áp trán vào cửa kính, ngắm khung cảnh đang lướt qua. ``Đây là\dots\ Aogami sao? Khác hẳn so với nơi em từng biết. Những ngôi nhà gỗ, những con đường lát đá\dots\ nhưng lại xen kẽ nhà bê tông, rồi cả xe điện nữa. Giống như một thế giới khác.''

Qua cửa kính, thị trấn trải dài dưới ánh sáng mùa đông. Aogami là một thị trấn nhỏ nằm giữa dãy núi xanh thẳm, nơi nhà gỗ truyền thống vẫn đứng cạnh tòa bê tông mới dựng. Đường lát đá uốn quanh sườn núi, dẫn xuống những cửa hàng tạp hóa bé nhỏ, quán cà phê ấm cúng. Dưới chân núi, những chiếc xe điện lăn bánh êm ru; trên cao, tiếng trống taiko từ ngôi đền cổ vọng xuống. Mọi thứ vẫn thanh bình, vẫn ``như vậy.'' Nhưng đồng thời\dots\ lại chẳng hoàn toàn là vậy. Giống một giấc mơ được tái tạo lại, từng chi tiết đúng nhưng tinh thần lại lạc đi đâu mất.

``Trông\dots\ giống hệt nơi chúng ta từng sống,'' Saikyou nói nhỏ, tay vuốt nhẹ cốc nước, ``nhưng đồng thời lại chẳng giống chút nào.''

``\dots Quả thật là như vậy. Anh cảm thấy nó có gì đó\dots\ lạc lõng hẳn,'' tôi nói, rồi liếc về người em gái kia. ``Cậu thì sao, Reiko-san?''

Reiko quay sang tôi, ánh mắt hai màu ánh lên vẻ kiên định, nét mặt trông có vẻ căng thẳng hơn, tỏ vẻ không vừa lòng. ``Onii-sama à.''

Tôi khựng lại. Tôi vừa nói gì sai ư?

Cô khẽ mỉm cười, rồi lắc đầu. ``Anh không cần xưng hô cứng nhắc như trong mơ đâu. Gọi em là `Reiko' thôi. Em cũng biết vì sao anh cứ gọi `cậu-tôi,' nhưng chúng ta là anh em, phải không?''

``Anh em\dots'' tôi lặp lại, nhíu mày. Thật sự, tôi vẫn chưa thể nào quen được cách gọi ấy.

Reiko gật đầu, giọng chắc nịch. ``Phải, anh và chị là người thân của em. Có thể ký ức bị xáo trộn, có thể mọi thứ không hoàn toàn rõ ràng, nhưng điều đó không thay đổi được sự thật này. Vậy nên\dots'' cô nở một nụ cười khẽ, một cách rất chân thành, ``xin anh\dots\ hãy đối xử với em như với Onee-sama, đừng coi em như là một người lạ nữa\dots\ được không ạ?''

Tôi mở miệng, nhưng không thốt ra lời nào. Ngực tôi thắt lại. Một phần trong tôi muốn phản bác: rằng đây chẳng qua là sự sắp đặt nào đó, rằng tất cả đều là giả tạo. Nhưng nhìn vào đôi mắt dị sắc đang dõi thẳng vào mình, ánh lên một niềm tin giản đơn, và một thứ gì đó thân quen mơ hồ trong lòng\dots\ tôi bỗng thấy bản thân chẳng có sức để phủ nhận.

``Reiko\dots'' tôi gọi, giọng run nhẹ. ``Vậy thì\dots\ được rồi. Tôi-- À không, anh cũng không ngại đâu.''

Nụ cười của cô nở rộng, ấm áp đến mức làm lòng tôi chao đảo. Saikyou cũng khẽ mỉm cười, như thể tán đồng với quyết định mà chính tôi còn chưa chắc chắn.

Thực tại này\dots\ tôi vẫn chưa hiểu. Nhưng ít nhất, lúc này đây, tôi có hai người em gái ngồi cạnh bên. Và điều đó cũng đủ để tôi tiếp tục ngồi yên trên chuyến tàu đang đưa chúng tôi đến ga Aogami.

Tôi chớp mắt một lần, hoàn toàn vô thức, và trong khoảnh khắc ngắn ngủi ấy\dots

``!?''

\dots cảnh tượng bên trong khoang tàu biến đổi hoàn toàn.

Những chiếc ghế xanh sậm trong toa tàu giờ đã phủ một lớp bụi dày đặc, mờ mịt như sương. Chúng bám chặt vào từng kẽ nứt trên bề mặt ghế, tạo thành những vệt xám xịt như thể đã tích tụ qua hàng thập kỷ không ai động đến.

Những tay nắm kim loại dọc lối đi treo vắt vẻo, rỉ sét loang lổ như những vết thương hoen ố. Chúng lung lay yếu ớt mỗi khi tàu rung lắc, phát ra những tiếng kẽo kẹt ghê rợn như tiếng rên rỉ của kim loại mục nát.

Trên những ô cửa kính, vô số dấu bàn tay đen sì chồng chéo, đan xen vào nhau như một mạng nhện khổng lồ. Những vết tay ấy in hằn sâu vào bề mặt kính, như thể ai đó - hay chính xác hơn là vô số bóng ma - đã tuyệt vọng cào cấu để thoát ra.

Và giữa khung cảnh hoang tàn ấy, những bộ bàn ghế trường học nằm ngổn ngang, vất vưởng khắp lối đi. Chúng xen kẽ với ghế ngồi của toa tàu, tạo thành một mê cung hỗn độn của đồ nội thất. Không khí vốn ấm áp giờ trở nên đặc quánh, lạnh buốt và ngột ngạt như thể có hàng trăm linh hồn vừa chen chúc quanh đây. Tại sao những bàn ghế trường học lại xuất hiện ở một nơi không tưởng như thế này?

\begin{figure}[H]
  \centering
  \begin{subfigure}[t]{0.45\textwidth}
    \centering
    \includegraphics[width=8cm]{../../../../Image/Book?/e2-1b1.png}
  \end{subfigure}
  \quad
  \begin{subfigure}[t]{0.45\textwidth}
    \centering
    \includegraphics[width=8cm]{../../../../Image/Book?/e2-1b2.png}
  \end{subfigure}
  \quad
  \begin{subfigure}[t]{0.45\textwidth}
    \centering
    \includegraphics[width=8cm]{../../../../Image/Book?/e2-1b3.png}
  \end{subfigure}
\end{figure}

Tôi nuốt khan, hai bàn tay siết chặt đầu gối trước cảnh tượng hoang tàn ấy.

``C-cái này\dots\ các em cũng thấy chứ?''

Không chỉ mình tôi, mà cả Reiko và Saikyou đều sững sờ với khung cảnh này.

Cô em song sinh liếc quanh, đôi mắt ánh lên sự cảnh giác: ``Ừ\dots\ không phải chỉ mình anh đâu, Nii-chan. Chỗ này\dots\ Quả thật có gì đó không ổn.''

Reiko khẽ run, nhưng lại cố bình tĩnh nhất có thể: ``Thế này thì\dots\ chắc chắn không phải trò đùa rồi. Onii-sama, em thấy khó thở quá\dots''

``B-Bình tĩnh nào, Reiko,'' tôi vội lên tiếng trấn an. ``Có lẽ đây chỉ là ảo giác thôi.''

``Nhưng em thấy\dots\ mọi thứ mà ta đang nhìn trước mặt\dots\ sao lại thực quá\dots'' Saikyou nói, mắt láo liên đảo quanh.

Cho dù chúng tôi có trấn tĩnh như thế nào, nhưng cả ba đều chỉ có thể ngầm hiểu rằng: khoang tàu này\dots\ hoặc có lẽ, là cả chuyến tàu này, không phải là một cái gì đó bình thường.

Không cần chờ đợi lâu để tìm lời giải đáp, một giọng nói trầm lặng vang lên từ phía đối diện: ``Đúng vậy, Kyuen-san, Saikyou-san, Reiko-san.''

Một giọng nói rất quen. Một giọng nói mà tôi đã nghe rất nhiều ở trong giấc mơ vừa mới chợp tắt ấy.

Chúng tôi cùng ngẩng đầu. Từ bóng tối, Misora Shino-san - cô gái với đôi mắt xanh sâu hút - lặng lẽ bước tới. Cô không mang theo vẻ đáng sợ như trong giấc mơ cũ, mà thay vào đó, ánh mắt cô tỏa ra một sự điềm tĩnh pha lẫn nghiêm trọng.

\img{e2-1c}

``Nhà ga Aogami vốn dĩ đã mang theo một hiện tượng siêu nhiên,'' cô ấy nói, giọng nhỏ nhưng vang vọng khắp khoang. ``Nó là hậu quả từ vòng lặp đã diễn ra\dots\ nơi ký ức, thời gian và bản thân con người bị cuốn trôi vào nơi này.''

Tôi nhíu mày. ``Ý cậu là\dots\ mọi thứ lặp đi lặp lại, và chúng tôi chỉ đang kẹt giữa nó?''

Cô gật khẽ: ``Các cậu chỉ mới bắt đầu cảm nhận được, và mọi thứ sẽ dần mất kiểm soát đấy. Nhưng\dots'' cô ngừng lại, ánh mắt quyết tâm hiện rõ trên gương mặt cô, ``đừng để những chuyện đó làm cho các cậu xao nhãng mà quên mất nhiệm vụ của chúng ta. Tìm và thuyết phục thực thể Sora quay trở lại. Cô ấy chính là mảnh ghép cần thiết để kết thúc chu kỳ này.''

Reiko cau mày: ``Còn cậu thì sao? Tại sao cậu không giúp trực tiếp?''

Đôi mắt Shino-san thoáng ánh buồn: ``Mình đã thử rất nhiều lần rồi, nhưng\dots\ như ta đã nói chuyện trước đó rồi đấy, mình đã không thành công, và mình đã mất hết sự kiên nhẫn rồi. Vậy nên\dots\ mình sẽ quan sát. Cách thức như thế nào là lựa chọn của các cậu.''

``Nói thì dễ lắm chứ\dots'' Saikyou gãi đầu, giọng có phần càu nhàu. ``Chưa kể đến chúng tôi còn chẳng biết cô ấy đang ở đâu nữa.''

``Mình tin các cậu sẽ làm được,'' Shino đáp, ánh mắt tỏ rõ sự nghiêm nghị. ``Các cậu cũng nên chuẩn bị tinh thần đi.''

Ngay sau câu nói ấy, thân ảnh của cô gái tóc nâu nhòa đi như một làn khói mỏng, để lại khoảng trống lạnh lẽo trước mặt chúng tôi, nhưng trước đó, cô ấy đã gửi cho tôi những lời động viên và chúc may mắn cho bọn tôi trước đó.

Có vẻ như mọi chuyện sẽ bắt đầu loạn kể từ đây.

Và cũng nhanh như lúc bắt đầu, khoang tàu trở lại bình thường: ghế ngồi sạch sẽ, cửa kính phản chiếu bóng núi mờ xa, tiếng bánh sắt nghiến trên đường ray rì rầm đều đặn.

Tôi thở phào, nhưng lòng ngổn ngang hơn bao giờ hết. Saikyou thì vẫn nghiêng đầu, trầm ngâm như cố xâu chuỗi những mảnh rời rạc vừa nghe được. Còn Reiko, vẫn ôm chặt cánh tay tôi, khẽ thì thầm: ``Em\dots\ vẫn không tin nổi là vừa rồi thật sự xảy ra.''

Tôi nhìn cả hai, lòng nặng trĩu, vì đến bản thân tôi cũng không thể ngờ được nó đã thực sự xảy ra, và cũng không chắc mọi chuyện từ đây sẽ diễn ra như thế nào thôi. Chỉ trong thoáng chốc, chuyến tàu này đã không còn đơn thuần là một chuyến đi đến nơi ở mới như ban đầu nữa.

Mỗi người chìm trong suy nghĩ riêng, nhưng cả ba đều im lặng, như thể sợ rằng chỉ cần cất tiếng, khoang tàu sẽ lại một lần nữa biến đổi.

Nó là khởi đầu cho một điều gì đó còn phức tạp hơn như vậy. Rất nhiều.

Ngoài kia, tiếng ray nghiến ken két hòa cùng tiếng máy móc đều đặn. Bên trong, khoang tàu vắng lặng đến lạ, chỉ lác đác vài bóng người ngồi xa tít. Tôi rút điện thoại từ túi quần, màn hình sáng lên. Đồng hồ chỉ 11:25. Chưa phải là giờ cao điểm. Mọi thứ dường như yên ổn\dots\ nếu bỏ qua những gì vừa xảy ra.

Đúng lúc ấy, một giọng nữ trong trẻo vang lên từ hàng ghế phía trước: ``Nào, Yuka, dậy đi! Chúng ta sắp đến ga Aogami rồi đấy!''

Kèm theo đó là tiếng than vãn lười biếng, kéo dài: ``Ưm\~{} Buồn ngủ quá\dots\ Khi nào tới thì tôi sẽ dậy mà, Akane\dots\ Tôi hứa đó\~{}''

Saikyou ngẩng lên, ánh mắt thoáng ngạc nhiên. Reiko cũng khựng lại, rồi nhìn sang tôi, như để chắc rằng không chỉ một mình em nghe thấy.

``Giọng nói ấy\dots\ quen thuộc ghê,'' tôi lẩm bẩm.

Khi chúng tôi tiến tới, khung cảnh dần rõ ràng: một cô gái tóc nâu, bờm tai chó cụp cùng màu, đang kiên nhẫn lay lay cô bạn có bờm tai mèo tím bên cạnh. Cảnh tượng vừa đời thường, vừa có gì đó\dots\ quen thuộc đến khó hiểu.

\img{e2-1d}

Cô gái tóc nâu thở dài, nhìn người bạn đang nằm co ro trên ghế với vẻ bất lực. ``Bà lúc nào cũng nói vậy, cuối cùng bà có chịu dậy đâu?'' cô lắc đầu, rồi nắm lấy vai bạn mình lắc nhẹ. ``Thôi nào, dậy đi! Không là tôi bỏ bà lại đấy!''

Cô gái tóc tím chỉ khẽ cựa mình, mắt vẫn nhắm nghiền. ``Ưm\~{} Rồi, rồi\dots\ tôi dậy ngay đây\dots'' cô lẩm bẩm, rồi lại cuộn tròn người, tiếp tục chìm vào giấc ngủ của mình.

``Trời ạ, thật tình!'' cô gái tóc nâu đưa tay lên trán, không giấu nổi vẻ chán nản trước thái độ lười biếng của bạn mình.

Hai anh em bất giác nhìn nhau. Trong đầu tôi, một cơn đau nhói nhẹ bùng lên, rồi hình ảnh vụt thoáng hiện: cô gái bờm chó này\dots\ đã từng đứng trước cổng trường, mắt đỏ hoe vô hồn, miệng lẩm bẩm từng chữ liên tục: ``Mau trả bạn cho tôi\dots\ Mau trả bạn cho tôi\dots''

Tôi siết chặt điện thoại, hít mạnh một hơi, cố phân biệt đâu là ký ức thật, đâu chỉ là ảo giác. Saikyou cũng khẽ nhăn mặt, đưa tay ôm trán, như thể em cũng vừa thấy điều tương tự.

``Onii-sama, Onee-sama\dots'' Reiko lên tiếng, giọng mang chút lo lắng. ``Hai người\dots\ làm sao thế?''

Tôi giật mình, vội vã đáp: ``K-Không có gì đâu.'' Em gái lớn cũng phụ họa, có phần gượng gạo: ``Ừ, chỉ là hơi mệt thôi.''

Nhưng cả hai chúng tôi đều không thể rời mắt khỏi cô gái bờm chó kia. Càng nhìn kỹ, càng có cảm giác quanh cô tỏa ra một thứ gì đó kỳ lạ, vừa quen vừa xa, vừa như một mối dây vô hình nối liền với chúng tôi. Ngay cả người bạn đang say ngủ bên cạnh cô - người mà có lẽ tôi chẳng chắc đã từng gặp hay chưa - cũng khiến lòng tôi bất an một cách mơ hồ.

Và rồi, như bị hút theo bản năng, ánh mắt của cô nàng tai chó từ từ quay lại phía chúng tôi. Tôi thoáng khựng, tim lại rộn nhịp, còn cô thì chỉ nghiêng đầu nhẹ, vẻ mặt như thắc mắc vì sao chúng tôi cứ nhìn chằm chằm vào cô ấy.

``Onii-sama, Onee-sama?'' Reiko kéo nhẹ tay áo chúng tôi. ``Hai người sao vậy?''

Tôi lắc đầu, cố giấu đi cơn hỗn loạn trong lòng: ``À\dots\ không có gì đâu. Chỉ là\dots\ anh thấy cô ấy giống như người nào đó đã từng gặp ở đâu rồi thôi.''

``Ở đâu cơ?'' Reiko hỏi thêm, nét mặt tò mò.

``T-Trùng hợp thôi. Không có gì đâu.'' tôi nói vội, kèm một nụ cười gượng. Bên cạnh, Saikyou khẽ gật đầu, ánh mắt đồng tình với lời tôi vừa dựng nên. Chúng tôi chẳng cần nói, nhưng cả hai đều hiểu rằng chuyện này nên giữ kín, ít nhất là lúc này.

\ct

Tiếng loa vang lên báo hiệu đoàn tàu sắp tới ga Aogami. Không khí trong toa khi đó bỗng chốc như bị hối thúc của thời gian. Tiếng xột xoạt của quần áo, tiếng kéo khóa vali hòa vào nhau tạo thành một bản hòa tấu của sự hối hả. Hành khách lục tục thu dọn đồ đạc, người thì kiểm tra lại túi xách, kẻ thì vươn vai sau chuyến đi dài. Ánh nắng gần trưa nhạt nhòa hắt qua cửa sổ, tạo những vệt sáng dài trên sàn tàu.

Còn ở phía trước, cô gái bờm chó vẫn đang ra sức lay gọi bạn mình, giọng nói của cô vang lên giữa không gian ồn ào của toa tàu.

``Nào, dậy đi chứ! Đến rồi còn gì, mau lên nào!''

Cô bạn tai mèo tím vẫn cuộn tròn, chỉ khẽ làu bàu: ``Ưm\~{} Buồn ngủ quá\dots\ Cho tôi ngủ thêm năm phút nữa đi\~{}''

``Trời ạ, cái bà này\dots!'' cô gái bờm chó nhăn mặt, rồi như chợt nghĩ ra, quay sang nhìn tôi một lần nữa.

``Cậu gì ơi\dots\ cậu có thể giúp mình gọi cô ấy dậy không?'' cô ấy giương đôi mắt nâu mang hàm ý nhờ vả. ``Không hiểu sao, tôi lay mãi mà chẳng nhúc nhích gì. Mọi khi cậu ấy có như vậy đâu.''

``T-Tôi á?''

``Ừm. Nhờ cậu đấy.''

Bị nhờ việc một cách đột ngột như vậy, tôi biết phải làm sao bây giờ? Tôi thoáng lúng túng, nhưng rồi vẫn gật đầu thực hiện: ``Ờ\dots\ được rồi.''

Tôi cúi xuống gần, gọi khẽ: ``Này\dots\ dậy đi, sắp đến ga rồi đấy.''

``Ưm\dots''

``Nếu cậu còn ngủ nữa thì tôi sẽ để bạn cô khiêng đi đấy.''

``C-Cái!?'' cô bạn bờm chó thốt lên.

``Onii-sama?'' em gái nhỏ của tôi cũng thoáng ngơ ngác với lời dọa dẫm có phần ngẫu hứng của tôi. Đừng có nhìn tôi vậy, còn cách nào khác để thuyết phục đâu!

``Ưm\dots\ Được rồi, được rồi\dots\ Dậy đây, dậy ngay đây\~{}''

Kỳ lạ thay, chỉ một tiếng gọi ấy, cô gái bờm mèo liền giật mình trở mình, rồi dụi mắt, ngáp một hơi dài như chưa từng ngủ say đến thế.

``Ưm\dots\ thật ư?''

``Thấy chưa!'' cô bờm chó bĩu môi, rồi khoác tay kéo bạn mình đứng dậy. Tôi và Saikyou chỉ kịp liếc nhau, trong lòng càng dấy lên thêm sự khó hiểu.

\ct

Khi đoàn tàu dừng hẳn, dòng người thưa thớt rời khỏi toa. Ba chúng tôi xuống theo lối hành lang ga. Phía trước, hai cô gái kia đi song song nhau: một người còn ngái ngủ, vừa đi vừa ngáp; một người thì loay hoay nhắc nhở bạn mình không được lề mề.

Bước chân tôi chậm lại. Tôi có cảm giác như những hình ảnh rời rạc trong đầu mình đang dần kết nối lại, và điều ấy khiến tim tôi bất giác đập nhanh hơn.

Ngay lúc đó, cô gái bờm chó khẽ ngoái đầu lại, ánh mắt hướng thẳng về phía chúng tôi.

``\hl{A-nô\dots{} c-cảm ơn cậu đã giúp\dots}'' cô gái bờm chó khẽ cúi đầu, giọng nhỏ nhẹ đến mức tôi suýt không nghe thấy. Khuôn mặt cô thoáng ửng hồng, như thể bản thân cũng không quen nói ra lời cảm ơn trực tiếp.

Tôi chớp mắt, bất giác mỉm cười gượng. ``À\dots\ không có gì đâu. Tôi chỉ nói đại thôi mà.''

Reiko nghiêng đầu, nhìn cả hai chúng tôi rồi khúc khích cười, nhưng không chen vào.

``V-Vậy à\dots'' cô khẽ lắp bắp, hai tay vân vê gấu áo. Cô gái bờm chó có vẻ không quen nói chuyện với người lạ, nhưng ánh mắt lại ánh lên niềm vui khi được trò chuyện.

Rồi, như vừa lấy hết dũng khí, cô ấy ngoảnh về chúng tôi, hai tay đan vào như thể hiện sự e thẹn. ``M-Mình là Inukai Akane. R-Rất hân nạnh được làm quen.''

\chrint

\begin{wrapfigure}[14]{l}{6cm}
  \includegraphics[height=8cm]{../../../../Image/Book?/akane_human.png}
\end{wrapfigure}

Inukai Akane (\ruby{戌}{いぬ}\ruby{飼}{かい}\ruby{茜}{あかね}, \ruby{I-nư}{Tuất}-\ruby{cai}{Tự}\ \ruby{A-ca-nê}{Thiến}), cô nàng bờm tai chó vừa mới bắt chuyện với nét mặt bối rối ấy.

Trước khi gặp Yuka, cô ấy hiếm khi thể hiện cảm xúc và cũng rất ít khi mỉm cười. Nhưng khi đã quen với ai đó, cô có thể đối xử bình đẳng và vui vẻ với họ, dù phải mất một thời gian để thực sự mở lòng. Tôi nghĩ cũng có một phần đúng, nếu chỉ xét riêng với những người bạn cùng giới, nhưng với những người con trai như tôi, đó có lẽ lại là một câu chuyện khác: cô ấy sẽ quay trở về với bản chất thật của mình, như chúng tôi đang thấy.

Cô luôn ngưỡng mộ cô bạn của mình vì cách sống thật với chính mình của cô ấy. Với Akane-san, không có gì hạnh phúc hơn là được đánh thức Yuka - người có thể chìm vào giấc ngủ ở bất cứ đâu. Đôi khi, cô tự hỏi mình sẽ ra sao nếu không còn Yuka nữa, bởi với cô, cô bạn thân ấy chính là tất cả. Cũng dễ hiểu thôi, vì cô ấy là người đã giúp cho Akane tự tin hơn khi giao tiếp với mọi người.

\chrint

``Akane nhút nhát lắm đấy\~{}'' cô gái bên cạnh - mà tôi nghĩ là bạn của cô ấy - ngáp dài, giọng lười biếng nhưng đầy ý tứ. ``Nhưng một khi đã quen thì cậu ấy vui vẻ lắm. Phải không nào?''

``Y-Yuka!'' Akane-san đỏ mặt, nhưng không phủ nhận. Cô liếc nhìn người bạn đang ngáp ngắn ngáp dài, rồi thở dài. ``Cậu lại buồn ngủ rồi à? Dù là nhà ngoại cảm nhưng cậu chẳng bao giờ để tâm mấy chuyện này nhỉ.''

``Ừm\~{} vì tớ thích sống thật với bản thân mà,'' Yuka cười khúc khích, mắt lim dim. ``Akane cũng nên thế đi\~{}''

\chrint

\begin{wrapfigure}[14]{r}{5cm}
  \includegraphics[height=8cm]{../../../../Image/Book?/yuka_human.png}
\end{wrapfigure}

Trái ngược với cô bạn nhút nhát nhưng quan tâm tới mình, Nekomiya Yuka (\ruby{猫}{ねこ}\ruby{宮}{いや}\ruby{友}{ゆ}\ruby{架}{か}, \ruby{Nê-cô}{Miêu}-\ruby{mi-gia}{Cung}\ \ruby{Giu}{Hữu}-\ruby{ca}{Giá}) lại có vẻ bình thản hơn. Cô ấy là người bạn thời thơ ấu của Akane, là một cô gái trầm tính với đôi mắt luôn híp lại như sắp ngủ, và thậm chí là ngủ khắp mọi nơi với một nụ cười hạnh phúc trên môi, đặc biệt là khi được cuộn mình dưới bàn kotatsu ấm áp vào mùa đông. Chẳng khác nào một con mèo nhà (mà trớ trêu thay, tên của cô cũng chứa cả Hán tự liên quan tới mèo.)

Dù sở hữu năng lực ngoại cảm mạnh mẽ có thể nhìn thấy những điều siêu nhiên mà không phải ai cũng cảm nhận được, cô dường như cũng chẳng bận tâm tới sức mạnh này lắm, mà vẫn hành xử như một cô gái bình thường. Cô đã quá quen với chuyện bị mọi người soi mới về năng lực này, và đôi lúc\dots\ chúng lại tỏ ra có ích.

\chrint

Xuống tới sảnh ga, dòng người tản dần, chỉ còn lại chúng tôi đi trên hành lang lát đá cũ. Cô gái bờm chó vẫn giữ khoảng cách nửa bước so với chúng tôi, như ngại ngần điều gì đó. Còn cô bạn mèo tím thì dụi mắt, ngáp thêm vài cái rồi mới cất tiếng, giọng có chút bông đùa: ``Akane có vẻ như thích nói chuyện với người lạ ghê\~{}''

``Thì\dots\ bởi vì chúng ta xuống cùng ga mà!'' Akane-san lúng túng đáp, tay khua khua trong không khí. ``Vậy nên tôi cũng\dots\ muốn mở rộng làm quen chứ!''

``Ha, tôi thì thấy bà lúc nào cũng co quắp lại, có chịu mở lòng ai bao giờ đâu\~{}'' Yuka-san lè lưỡi trêu chọc, khiến cô bạn tóc nâu đỏ bừng mặt.

``Y-Yuka!''

Tôi bật cười khẽ, khẽ nghiêng người sang Reiko: ``Quả thật cô ấy thú vị nhỉ. Nhìn thì hiền hòa, thân thiện, nhưng đúng là hơi nhút nhát.''

``Đúng thật nha\~{}'' Reiko gật đầu đồng tình, đôi mắt ánh lên vẻ thích thú, như thể vừa tìm thấy điều gì mới mẻ.

Chúng tôi rời khỏi toa tàu, bước vào sảnh ga Aogami. Không gian rộng lớn của nhà ga vang vọng tiếng bước chân và tiếng trò chuyện rì rầm của hành khách. Ánh sáng từ những ô cửa kính cao chiếu xuống, tạo những vệt sáng dài trên nền gạch bóng loáng. Mùi gỗ thông từ những thanh xà ngang trên trần nhà hòa quyện với hương cà phê thoang thoảng từ quầy bán hàng nhỏ góc sảnh.

``À mà\dots'' Akane-san chợt chuyển chủ đề, có lẽ muốn trốn khỏi màn chọc ghẹo của cô bạn. ``Điều gì đã đưa các cậu tới Aogami vậy?''

``Ừ nhỉ. Mình cũng thấy tò mò đấy,'' Yuka-san nối thêm, giọng vẫn còn ngái ngủ nhưng đầy hiếu kỳ.

\img{e2-1e}

Tôi và Saikyou thoáng nhìn nhau. Một cái gật, một cái nhíu mày, chỉ bằng cử chỉ, cả hai đã ngầm thống nhất sẽ né tránh. Nhưng ngay khi tôi định mở lời, Reiko đã tươi cười chen ngang: ``Bọn mình tới nhập học vào mùa hè này, nên ba bọn mình muốn tham quan thị trấn này một vòng để dần thích nghi với nơi ở mới hơn!''

Cả hai chúng tôi đồng loạt giật mình: ``Này, Reiko--''

``Ơ? Chẳng phải chúng ta tới đây để làm việc đó sao?'' em ấy chớp mắt, nhưng rồi ánh nhìn của em khẽ chùng xuống, như đã nhận ra hai chúng tôi đang che giấu điều gì đó.

Nhưng, thay vì tra hỏi, em ấy chỉ gật gù như thể tự xác nhận, rồi tiếp tục dẫn dắt cuộc hội thoại. Dù gì, đây cũng là lần đầu - theo đúng nghĩa đen - Reiko gặp hai người đó, nên để em út dẫn dắt có lẽ sẽ tự nhiên và dễ nói chuyện hơn.

``Vậy các cậu đã có chỗ ở chưa?'' Yuka-san nghiêng đầu hỏi.

``Nhà mới của bọn mình cũng ở gần đây thôi,'' em út đáp nhanh nhẹn, vẫn với dáng vẻ tự tin và hòa hảo với mọi người như ở trong giấc mơ - hoặc có lẽ, từ vòng lặp trước, ``Mình nghĩ là đâu đó cách trường khoảng mười lăm phút đi bộ đấy.''

``May nha\~{}'' cô bạn mèo vươn vai, cười híp mắt. ``Sắp tới chúng ta sẽ là hàng xóm của nhau rồi hen\~{}''

``Bọn mình cũng muốn các cậu dẫn dắt tham quan vòng quanh thị trấn này,'' Akane-san nói thêm, lần này giọng có vẻ chắc chắn hơn.

``Ờm, bọn mình có thể tự đi được--'' tôi định mở lời từ chối, nhưng cô ấy lập tức lắc đầu, đôi mắt ánh lên sự cứng rắn hiếm thấy.

``Không phải khách sáo! Bọn mình tự nguyện làm việc này cho các cậu mà!''

Tôi nuốt khan. Trước ánh nhìn đó, tôi không thể nào nói ``không'' thêm lần nữa. ``\dots Được thôi.'' Chỉ chờ có thế, cô nàng tỏ vẻ vui ra mặt, toàn thân giãn ra.

Yuka-san mỉm cười, tỏ vẻ đồng cảm với tôi. ``Akane một khi đã quyết tâm làm gì đó thì sẽ làm tới cùng. Đó là điểm mà mình quý cậu ta đấy\~{}''

``Vậy\dots\ chín giờ sáng mai nhé?'' Akane liếc chiếc điện thoại, mở lịch biểu ra và gõ gõ vài thứ trên đó. ``Ta sẽ gặp nhau tại khu vực ở khu phố truyền thống, cũng quanh đó thôi.''

``Ừm, được rồi,'' Saikyou đáp. ``Cũng không xa lắm.''

Chúng tôi bước tiếp trên hành lang ga, tiếng bước chân vang vọng dưới trần mái vòm. Tôi liếc sang em gái lớn, thấy em cũng đang lặng lẽ quan sát cô bạn tai chó. Trong lòng tôi dấy lên một nỗi chắc chắn: cô gái ấy chính là người đã đứng ngoài cổng trường, lẩm bẩm câu ``Mau trả bạn cho tôi'' lặp đi lặp lại.

Điều mà tôi không thể hiểu được\dots\ là vì sao cô lại ở đó. Và\dots\ tại sao cô ấy đứng đó, lặp đi lặp lại đúng một câu như vậy? Điều duy nhất mà tôi có thể liên hệ được là nó chắc chắn liên quan tới vòng lặp quái quỷ kia. Nhưng phán đoán là một chuyện, chỉ ra để chứng minh lại là một câu chuyện khác.

\dots Thôi thì, tạm gác bỏ mấy thứ siêu nhiên đó ra một bên, có lẽ tôi và hai em gái của tôi sẽ phải làm quen với môi trường vừa mới mẻ, vừa quen thuộc này.

\ct

Khi chúng tôi bước ra khỏi cửa chính, không khí lạnh của tháng Hai ùa vào mặt. Trước mắt chúng tôi là quảng trường nhà ga - một khoảng không gian rộng lớn lát đá granite xám, với đài phun nước ở trung tâm đang ngủ đông. Xa xa, những dãy núi xanh thẫm ôm lấy thị trấn như một vòng tay khổng lồ. Những mái nhà truyền thống xen lẫn với công trình hiện đại tạo nên một bức tranh hài hòa giữa những gì xưa cũ và những gì mới mẻ, như một bản điêu khắc nghệ thuật thường thấy ở các khu phố cổ.

Ngay khi chúng tôi vừa bước tới cửa ga, một giọng con gái chợt vang lên, có chút trầm mặc, có chút khàn: ``À-rế? Kia chẳng phải là Akane-chan với Yuka-chan đó sao?''

Cả ba chúng tôi - tôi, Saikyou và Reiko - ngẩng đầu về phía phát ra giọng nói ấy. Trước mắt chúng tôi là ba cô gái khác, cũng mặc đồng phục giống hệt, như thể đã hẹn trước để chờ ở đây.

\img{e2-1f}

Người đầu tiên thu hút ánh nhìn của tôi là một cô gái tóc đen cắt ngắn, mắt xanh biếc, tay cầm hộp sữa con vẫn còn dở dang. Cách cô mỉm cười tự tin và nâng hộp sữa lên trông đầy năng lượng, dễ tính đến mức khiến tôi thấy thoải mái ngay từ cái nhìn đầu tiên.

Kế bên là một cô gái tóc vàng, mái tóc sáng óng ánh dưới nắng ban trưa, có cài thêm chiếc kẹp hình con dơi đính hạt ngọc màu cam. Giọng nói của cô khi cất lên thì ngân dài, như thể cố tình kéo lê từng âm, tạo ra cảm giác dịu dàng và lười biếng một cách dễ chịu.

Và cuối cùng là một cô gái có bờm tai sói nhô lên, mái tóc đen xen vài vệt đỏ nổi bật. Điều khiến tôi chú ý không chỉ là ánh mắt điềm tĩnh, mà còn là bộ đồng phục được mặc gọn gàng chỉnh tề, khác hẳn sự xuề xòa của hai người bạn cạnh mình. Trưởng thành, nghiêm túc, tôi đã nghĩ ngay như thế.

``Ái chà? Kaoru-san với mọi người cũng ở đó kìa,'' Yuka-san vươn vai, giờ thì tỉnh táo hẳn, vẫy tay chào.

``Chào buổi trưa mấy cậu!'' cô gái tóc ngắn giơ hộp sữa lên, vừa cười vừa nói. ``Và\dots\ dường như tôi thấy hai cậu vừa kết thân với ba người mới nhỉ?''

``Đúng là vậy đó! Bọn mình vừa mới làm quen với họ xong!'' Akane-san nhanh nhảu đáp, hơi ấp úng nhưng đầy phấn khởi. ``Mình còn hứa với ba người họ là ngày mai sẽ dẫn đi tham quan thị trấn này nữa!''

``Quả thật là Akane-chan vẫn luôn muốn kết thân hẳn với người mới nhỉ,'' Cô gái tóc ngắn gật gù, vẫn giữ nụ cười dễ gần.

``À, tiện thể thì xin giới thiệu,'' Cô liếc về phía chúng tôi rồi đưa hộp sữa sang tay còn lại, đưa tay kia ra một cách hoạt bát. ``Mình là Ukai Kaoru.''

\chrint

\begin{wrapfigure}[14]{l}{7cm}
  \includegraphics[height=8cm]{../../../../Image/Book?/kaoru.png}
\end{wrapfigure}

Ukai Kaoru (\ruby{鵜}{う}\ruby{飼}{かい}\ruby{薫}{かおる}, \ruby{U}{Đề}-\ruby{ca-i}{Tự}\ \ruby{Ca-ô-rư}{Huân}). Một trong những thủ lĩnh được yêu mến nhất trong lớp, luôn tỏa ra năng lượng tích cực và sẵn sàng giúp đỡ mọi người, chẳng mấy để ý đến khoảng cách ban đầu giữa người lạ với nhau.

Dù thường tỏ ra ngây thơ và dễ tính trong cuộc sống hàng ngày, nhưng khi đối mặt với những tình huống quan trọng, cô lại thể hiện sự điềm tĩnh đáng ngạc nhiên. Có lẽ vì vậy mà bạn bè thường tìm đến cô để xin lời khuyên, và họ luôn tin tưởng vào những gì cô nói.

Tuy nhiên, cô có một nỗi bận tâm nhỏ - mà nghe qua có thể khiến người ta bật cười: chiều cao của mình. Vì thế, ngày nào cô cũng uống sữa, hy vọng có thể cao thêm được một chút. Nhưng điều đáng quý nhất ở cô có lẽ là tấm lòng nhân hậu - khi thấy ai đó gặp khó khăn, cơ thể cô thường phản ứng nhanh hơn cả suy nghĩ để giúp đỡ họ.

\chrint

Ngay khi Kaoru-san vừa dứt lời, cô gái tóc vàng cũng bước lên, nở nụ cười dịu dàng: ``Này, này, ba người họ đi cùng các cậu là ai vậy?'' Giọng nói cô ấy kéo dài, từng chữ như trôi ra chậm rãi.

Tôi nhìn sang Saikyou, rồi cả ba chúng tôi tự giới thiệu. ``Tôi là Haragen Kyuen, còn đây là Saikyou, bọn tôi là anh em sinh đôi. Còn đây là Reiko, em gái của chúng tôi.''

``Hể\~{} Bảo sao mà hai người giống y đúc ghê\~{}'' co nàng tóc vàng đáp lại, rồi liếc về phía em gái út. ``Reiko-chan trông cũng đáng yêu thật đấy\~{}, đôi mắt hai màu của em trông cũng dị ghê\~{}''

Reiko mỉm cười, đáp lại: ``Mình không dám tự nhận bản thân dễ thương đâu. Cậu cứ nói quá. À mà, cậu là ai vậy?''

``Rất hân hạnh được làm quen nhé\~{}''' cô gái tóc vàng chắp tay nhẹ nhàng. ``Mình là Kisaragi Saya.''

\chrint

\begin{wrapfigure}[14]{r}{4cm}
  \includegraphics[height=8cm]{../../../../Image/Book?/saya.png}
\end{wrapfigure}

Kisaragi Saya (\ruby{如月}{きさらぎ}\ruby{星}{さ}\ruby{夜}{や}, \ruby{Ki-xa-ra-ghi}{Như Nguyệt}\ \ruby{Xa}{Tinh}-\ruby{gia}{Dạ}), cô gái với chất giọng dịu dàng và trấm ấm như một người mẹ.

Tôi không khỏi chú ý đến cách nói chuyện đặc biệt của cô ấy - luôn kéo dài những âm cuối câu một cách từ tốn, như thể đang đắm chìm trong suy nghĩ của riêng mình. Đó là một trong những nét đặc trưng khiến Saya-san trở nên dễ gần trong mắt mọi người.

Sở trường của cô ấy là các vì sao trên trời. Một khi Saya-san bắt đầu nói về chúng, cô ấy có thể say sưa hàng giờ liền. Đó là niềm đam mê cô ấy ấp ủ từ nhỏ - dành cả ngày để ngắm bầu trời và mơ mộng về những thiên thể xa xôi. Nhờ vậy mà kiến thức thiên văn của cô ấy phong phú đến đáng ngạc nhiên.

\chrint

Akane-san lúc này chen ngang, phấn khích như thể tiết lộ bí mật: ``Này nhé, biết gì không? Cả ba người họ sẽ chuyển tới trường Aogami chúng ta vào mùa hè này đấy!''

``Ồ, thật sao?'' người có bờm tai sói cuối cùng mới cất tiếng, giọng điềm đạm hơn hẳn hai người bạn. Đôi mắt cô nhìn thẳng vào chúng tôi, như muốn cân nhắc kỹ càng từng lời. ``Nhưng nếu vậy thì\dots\ tại sao họ lại ở đây?''

Cô bạn bờm mèo khựng lại một thoáng, mắt lướt sang Yuka, rồi ngập ngừng: ``À, chuyện là thế này\dots''

Reiko nhanh nhảu tiếp lời, giọng lanh lợi: ``Bọn mình được ba dẫn tới đây để tham quan thị trấn, để làm quen trước khi nhập học! Thực ra bọn mình cũng mới vừa đặt chân tới thôi.''

Cả Akane và Yuka gật gù phụ họa, thậm chí Yuka còn chống nạnh ra vẻ xác nhận: ``Chuẩn luôn đấy. Đúng như thế luôn.''

Tôi và Saikyou chỉ có thể im lặng, xen vào vài câu ``Ừ'' và ``Đúng vậy'' cho khớp, để cô em út thoải mái điều khiển nhịp câu chuyện. Không rõ từ bao giờ, cô em gái này đã quen với việc đưa chúng tôi thoát khỏi những tình huống khó xử như thế.

Ngay khi Reiko dứt lời, Kaoru-san là người hăng hái nhất, cô nhấc hộp sữa lên, nghiêng đầu một cái rồi bật cười: ``Thế thì tốt quá rồi! Nếu ngày mai Akane-chan với Yuka-chan định dẫn mấy cậu đi tham quan thị trấn, sao không bàn kỹ luôn nhỉ? Tôi cũng muốn nghe kế hoạch của hai người xem sao.''

``Đúng đó\dots' Akane-san đáp, khoanh tay lại suy nghĩ. ``Mình đang tính là sẽ dẫn họ ra trung tâm thị trấn. Ở đó có nhiều cửa hàng, với cả quảng trường nữa. Sau đó có thể tới đền thờ trên đồi, vì mùa này không đông người lắm.''

Đền thờ? Nghe có vẻ quen quen\dots\ Hình như tôi đã từng thấy nó trong giấc mơ thì phải\dots

Tôi liếc sang cô em gái lớn, nhìn thấy mặt em ấy cũng khẽ nhíu mày. Tôi có thể đoán em ấy đang nghĩ: ``Tại sao Akane-san lại nhắc đến ngôi đền nhỉ\dots''

Bình thường, việc nhắc tới ngôi đền chẳng có gì là lạ cả. Nó là một phần không thể thiếu của thị trấn Aogami, một điểm đến quen thuộc của cả người dân lẫn du khách. Nhưng lần này thì khác. Tôi không biết tại sao, nhưng cứ mỗi lần ai đó nhắc tới ngôi đền, tôi lại cảm thấy có gì đó không ổn. Một cảm giác khó tả, như thể có một thứ gì đó đang chờ đợi chúng tôi ở đó.

Tôi liếc nhìn Saikyou, rồi cả Reiko. Cả hai đều có phản ứng tương tự - em gái sinh đôi của tôi thì im lặng đầy ý nghĩa, còn Reiko thì khẽ giật mình, dù cố che giấu. Có vẻ như không chỉ mình tôi cảm nhận được điều bất thường này.

Yuka-san vươn vai, cười khúc khích: ``Rồi tiếp theo thì\dots\ ăn vặt thôi! Tôi đảm bảo quán bánh cá gần ga là số một đó nha\~{}''

Saya-san vỗ tay khe khẽ, giọng ngân dài: ``Nghe hay ghê\~{} Mình cũng muốn đi theo lắm luôn đó\~{}''

``Ừm,'' cô gái tai sói - trời đất, tôi còn chưa biết tên của cô ấy - chắp tay sau lưng, bước đi chậm rãi hơn, ánh mắt liếc qua tôi rồi sang Saikyou. ``Kế hoạch sơ qua khá hợp lý. Có lẽ nếu cần thì chúng ta cũng có thể chia nhóm nhỏ, đề phòng lạc nhau.''

Reiko thì mắt sáng rực, gần như reo lên: ``Nghe thú vị quá! Chưa gì mà mình đã mong chờ buổi đi tham quan ngày mai rồi đấy!''

Tôi chợt nhận ra bản thân vẫn còn lộn xộn với các tên gọi. Tôi giơ tay, bắt đầu chỉ từng người, vừa nói vừa xác nhận lại từng người một trong số năm cô gái kia.

``Được rồi, xem nào\dots\ giờ ta có Akane-san\dots'' Tôi chỉ sang cô gái bờm chó đang vui vẻ.

``\dots Yuka-san\dots'' người bạn bờm mèo đang gãi đầu ngái ngủ.

``\dots Kaoru-san\dots'' cô gái tóc ngắn hớn hở giơ hộp sữa đáp lại.

``\dots Saya-san\dots'' cô gái tóc vàng khẽ cúi đầu, nụ cười dịu dàng vẫn nở.

``Và\dots\ Ờm\dots'' Tôi khựng lại một nhịp trước cô gái bờm sói, chưa kịp mở miệng. ``Cậu tên là gì vậy?''

``Nhắc mới nhớ, hình như Miku-chan chưa giới thiệu bản thân nhỉ?'' Kaoru-san lên tiếng.

Cô liền cười nhẹ, lên tiếng: ``À, xin thứ lỗi. Từ nãy chúng ta mải nói chuyện mà tôi lại quên mất. Tôi là Kamishiro Miku. Mong được làm quen.''

\chrint

\begin{wrapfigure}[14]{l}{8cm}
  \includegraphics[height=8cm]{../../../../Image/Book?/miku_human.png}
\end{wrapfigure}

Kamishiro Miku (\ruby{神}{かみ}\ruby{代}{しろ}\ruby{美}{み}\ruby{玖}{く}, \ruby{Ca-mi}{Thần}-\ruby{si-rô}{Đại}\ \ruby{Mi}{Mỹ}-\ruby{cư}{Cửu}) là một cô gái thanh tao và chín chắn hơn so với các bạn cùng trang lứa. Theo lời kể của Kaoru-san và Saya-san, cô ấy luôn cố gắng tỏ ra trưởng thành và thanh lịch, đến mức sưu tập đầy đủ các tạp chí thời trang phụ nữ. Tuy nhiên, đằng sau vẻ ngoài sang trọng ấy lại là một tâm hồn dễ xúc động - cô ấy thường đỏ mặt ngượng ngùng mỗi khi được ai đó khen ngợi.

Điều thú vị là dù cô ấy luôn cố gắng nói năng như một tiểu thư con nhà giàu có, nhưng khi về nhà, giọng điệu đó lại biến mất hoàn toàn, thay vào đó là một Miku-san đáng yêu và tự nhiên hơn nhiều. Theo như lời họ, cô ấy bắt đầu ăn nhiều sô cô la, chỉ vì đọc được một bài báo nói rằng nó tốt cho làn da. Quả thật là một cô gái có những điểm đối lập thú vị.

\chrint

``Ừm\dots\ vậy là đủ rồi,'' tôi khẽ gật đầu, lòng có chút an tâm hơn.

Cả năm cô gái bắt đầu quay sang hỏi thêm về bản thân chúng tôi. Reiko hăng hái kể trước, còn tôi và Saikyou thì đành gật gù phụ họa, thêm đôi ba chi tiết như là lấy lệ. Câu chuyện cứ thế xoay vòng: Kaoru-san hỏi về môn thể thao ở trường, Saya-san tò mò về sở thích riêng, Yuka-san cười nói vu vơ về giấc ngủ trưa, Akane-san thì bắt đầu tự tin góp chuyện hơn ban nãy, còn Miku-san thì giữ vai trò cân bằng, thỉnh thoảng chen vào những câu gọn gàng, súc tích.

Tôi không nhớ hết ai nói gì, chỉ biết rằng suốt gần một tiếng, tám người chúng tôi cùng nhau vừa đi vừa trò chuyện như thể đã thân quen từ lâu.

Ngoài kia, thị trấn Aogami trải dài trong bầu không khí mùa đông nhè nhẹ. Hơi thở phả ra khói trắng, những hàng cây ven đường khẳng khiu, vài chiếc lá đỏ cuối mùa vẫn cố bám trụ. Tiệm bánh cũ kỹ tỏa hương thơm ngọt, còn từ xa, tiếng chuông gió leng keng từ mái hiên một cửa tiệm nhỏ vọng lại. Đường phố không quá đông, nhưng từng ánh đèn cửa hàng, từng dòng người lác đác khiến khung cảnh vừa yên bình vừa sống động.

Chúng tôi đi trong khung cảnh ấy. Ttám người, tám nhịp bước. Tiếng cười nói hòa lẫn vào hơi gió lạnh, để lại trên con phố một dấu ấn kỳ lạ mà tôi biết, khó lòng quên được.

Thôi thì, cứ tận hưởng khoảng thời gian yên bình hiếm hoi này thôi nhỉ?

\pov{Haragen Saikyou}

Khoảng một giờ chiều.

Mặt trời đã lên tới đỉnh, ánh nắng buổi trưa dội xuống những con đường lát gạch một cách sạch sẽ. Cái bụng no căng sau bữa ăn khiến bước chân tôi có phần nặng nề, nhưng nhờ làn gió thoảng qua, cả nhóm vẫn thong dong đi về phía công viên để nghỉ ngơi.

Chúng tôi đi theo hàng lộn xộn: Kaoru-san bước phía trước, hộp sữa vẫn còn lắc lư trong tay (tôi không rõ cô ấy uống chậm, hay có lẽ cô mang cả lốc); Akane-san và Yuka-san tụm lại một góc trò chuyện nho nhỏ; còn tôi, Nii-chan và Reiko đi gần nhau hơn, giữ một khoảng cách vừa đủ.

``Vậy là, chúng ta đã có một kế hoạch sơ sơ cho sáng mai rồi,'' cô bạn tóc ngắn bất ngờ lên tiếng, như thể muốn chốt lại cuộc trò chuyện trước bữa trưa.

Miku-san vừa bẻ một miếng sô-cô-la từ thanh mua ở cửa hàng tiện lợi, vừa nhìn chúng tôi: ``Kế hoạch của Akane-san và Yuka-san thì thế, nhưng ý kiến của các cậu cũng quan trọng chứ. Có địa điểm nào cụ thể mà các cậu muốn tham quan không?''

Tôi đưa tay vuốt nhẹ mái tóc đuôi ngựa, nghĩ một lát rồi đáp: ``Mình muốn tìm một nơi nào đó sôi động trong thị trấn\dots\ chứ không phải chỗ yên tĩnh quá.''

Ngay sau đó, Reiko thêm vào: ``Mình thì muốn một nơi mà mọi người đều có thể đi cùng nhau. Miễn sao ai cũng thấy vui.''

Cả hai chúng tôi đều hướng ánh mắt về phía Nii-chan, chờ anh ấy đưa ra một ý kiến chắc nịch để làm điểm tựa. Thế nhưng, anh chỉ thốt ra một câu nhẹ bẫng: ``Đâu cũng được.''

Cả hai chị em tôi đồng loạt đổ mồ hôi hột, chẳng biết phải tiếp lời thế nào. Cũng không ngạc nhiên lắm với ông anh dễ tính tới độ dễ làm người ta nổi cáu đâu.

Chính lúc đó, Saya-san xen vào bằng một nụ cười thoáng qua: ``Nếu vậy thì các cậu thử đến ngôi đền Aogami xem nào?''

\img{e2-1g}

Ba chúng tôi thoáng khựng lại. Tôi nhíu mày, Reiko cũng nghiêng đầu khó hiểu, còn Nii-chan thì có đôi chút ngạc nhiên. Thật khó mà không để ý: chẳng phải vừa mới đó, Akane-san cũng đề cập đến ngôi đền này sao? Và rõ ràng, điều đó chẳng liên quan gì tới ``sôi động'' hay ``mọi người có thể cùng đi'' cả.

Nhưng trước khi tôi kịp mở miệng, Kaoru-san đã lập tức chen ngang, giọng bình thản: ``Đó cũng là tên của ngôi đền vùng này.'' Hộp sữa trên tay cậu ta khẽ phát ra tiếng lách tách khi cậu ta lắc nhẹ.

``R-Ra vậy\dots'' tôi gượng cười, dù chẳng ai hỏi gì mình. Một nỗi bất an mơ hồ dấy lên trong lòng. Từ khi cái tên ``ngôi đền'' được nhắc đến, tôi đã cảm thấy có gì đó mờ ám, như thể có thứ đang che giấu sau lớp vỏ quen thuộc ấy.

Tôi quyết định lái câu chuyện sang hướng khác, để bớt đi cảm giác là lạ khi cái tên ``ngôi đền'' cứ ám tới chúng tôi như thế.

``À\dots\ này,'' tôi cất tiếng, ``năm người các cậu có học chung trường không?''

Vừa hỏi xong, tôi mới chợt nhận ra sự ngớ ngẩn trong câu hỏi đó. Cả năm cô gái đều mặc chung một loại đồng phục, ít nhất là những gì tôi thấy ở chân váy xanh và chiếc nơ thắt ngực đồng màu, giống như tôi và Reiko đang diện. Nhưng thôi, đã lỡ nói thì chẳng rút lại được nữa.

Kaoru bật cười khẽ, như thể câu hỏi của tôi thú vị lắm vậy: ``Tất nhiên rồi! Đã thế nha, còn chung lớp nữa chứ.''

Câu trả lời ấy bỗng khiến tôi giật mình. Có một cái gì đó rất quen thuộc, một công thức lặp lại mà tôi và Nii-chan của tôi từng trải qua. Ở vòng lặp trước, vào tháng Bảy năm 20XX, chúng tôi đã học cùng một lớp với năm người: Suzune, Saki, Touka, Inari và Hanabi.

Rồi khi trôi ngược về quá khứ năm 1946, ngoài Shino và Reiko ra, lại là năm người bạn khác: Nanase, Yae, Honoka, Sakura và Shiina.

Và điểm chung? Tất cả đều học cùng một lớp.

Kaoru đột nhiên hạ giọng, thoáng chút ngập ngừng: ``Cơ mà gần đây ấy\dots\ rất nhiều bạn học trong lớp tôi đã không tới trường rồi. Có lẽ trận cúm mùa lại khởi phát. Cũng sắp tới xuân rồi mà. Đã thế, chẳng còn ai liên lạc được với họ nữa.''

Tôi sững người. Một mảnh ký ức khác lại nổi lên. Ở vòng lặp trước, tôi cũng từng nghe một chuyện y hệt: bạn học dần dần không đến lớp, và trận cúm mùa xảy tới. Tim tôi chùng xuống.

``V-Vậy à\dots'' tôi buột miệng đáp, cố giữ giọng bình tĩnh, nhưng đầu óc thì quay cuồng.

Một cơn déjà vu khác ập tới, mạnh mẽ hơn hẳn. Touka và Inari đã từng nhắc đến ``trận cúm mùa'' đó; nó không đơn thuần như vậy. Và ngay lập tức, ký ức kinh hoàng ùa về: trận động đất rung chuyển mặt đất, rồi trận sạt lở vùi lấp mọi thứ. Tôi và Nii-chan đã tận mắt chứng kiến những bạn học bỏ mạng trước mắt mình, và bản thân chúng tôi cũng đã suýt nằm lại sau những sự kiện khó quên ấy.

Tôi rùng mình. Một cơn gió lạnh thổi đến làm tôi run lên bần bật, cảm giác như vừa thổi vào sống lưng.

``Onee-sama? Chị ổn chứ?'' giọng Reiko vang lên đầy lo lắng.

``Cậu có ổn không vậy? Mặt cậu tái mét kia kìa\dots\ Cảm thấy không khỏe à?'' Kaoru-san cũng quay sang hỏi, lông mày nhíu lại.

Tôi hít một hơi sâu, lắc đầu và cố mỉm cười gượng gạo:
``T-Tôi ổn mà. Chỉ là hơi mệt thôi.''

Dù biết chỉ là một lời lấp liếm, tôi chẳng muốn làm ai phải lo thêm nữa. Nhưng trong lòng, tôi lại càng thấy bất an hơn bao giờ hết. Những sự trùng hợp này\dots\ hẳn phải có một gì đó liên quan tới một cái gì đó mà tôi không thể tìm ra được ngay lúc này.

Tôi liếc về phía Nii-chan. Anh ấy cũng đang toát mồ hôi hột, khuôn mặt căng cứng chẳng kém gì tôi sau khi nghe Kaoru-san nhắc tới chuyện ấy. Tôi chỉ khẽ gọi, thật nhỏ như sợ không để họ nghe thấy.

``\hl{Nii-chan\dots}''

Anh không đáp, nhưng ánh mắt nghiêm trọng kia đã nói thay lời. Cái gật đầu chậm rãi của cậu khiến cả lòng tôi như thắt lại. Không cần nói một câu gì thêm, chỉ cái cử chỉ ấy thôi cũng để cho hai em tôi có thể hiểu nhau.

Đúng lúc đó, cô nàng Miêu cất giọng. Lạ lắm, không còn cái kiểu lè nhè lười nhác nữa, mà khô khốc, nghe như cào vào tai: ``Nếu mà không ai liên lạc được với họ thì\dots\ quả đúng là kỳ lạ thật đấy.''

Kaoru-san gật đầu đồng tình, gương mặt cũng nghiêm lại.

Rồi, Yuka nói tiếp, chậm rãi, như từng chữ đang trượt ra từ khoảng tối nào đó: ``Mà nhắc mới nhớ, dạo gần đây \emph{càng có nhiều người mất tích ở thị trấn Aogami} này\dots''

Kaoru-san lập tức tiếp lời: ``Nghe người ta đồn, \emph{lần cuối cùng mà những người đó liên lạc được\dots\ là ở gần khu vực ga Aogami.}''

\img{e2-1h}

Tôi nín thở. Ngay khoảnh khắc ấy, cả ba anh em chúng tôi đồng loạt nhớ lại lời Shino-san đã dặn trên chuyến tàu lúc sáng, trước khi chính thức bước vào vòng lặp này.

``Tàu đến ga Aogami vốn dĩ đã mang theo một hiện tượng siêu nhiên. Nó là hậu quả từ vòng lặp đã diễn ra\dots\ nơi ký ức, thời gian và bản thân con người bị cuốn trôi.''

Đúng thật. Mọi chuyện đang trượt dần vào quỹ đạo kỳ quặc đó.

Tôi thì thầm với hai người kia, giọng run run: ``Nii-chan\dots\ Reiko\dots\ Có phải\dots\ chuyện Shino-san nói, đang bắt đầu hiện ra rồi không?''

Em tôi ôm cánh tay mình, khẽ rùng mình: ``Em\dots\ em cũng nghĩ vậy. Nghe cứ như là--''

Kyuen cắt lời, thấp giọng:
``Không. Nó không chỉ `cứ như' đâu. Đây chính xác là thứ cô ấy muốn cảnh báo chúng ta.''

Chúng tôi nhìn nhau, ba anh em đều chung một nỗi sợ hãi, nhưng lại chẳng ai dám nói to hơn.

Bỗng Yuka trầm xuống, giọng cô ma mị hẳn, nghe như vọng ra từ hang tối: ``Nhưng\dots\ nếu những bạn học trong lớp mình ấy\dots\ họ không hề bị cúm\dots\ mà lại là một phần trong số những người mất tích kia thì sao?''

Không khí nặng như chì. Một áp lực vô hình đè lên lồng ngực tôi. Tôi lập tức liên hệ tới lời Shino-san, và sự thật ấy càng trở nên lạnh lẽo, trơ trọi.

Thế rồi, đột ngột\dots

``\emph{Chỉ là cúm mùa thôi.}''

Một giọng nói lạ vang lên. Nó không phát ra từ bất cứ ai trong tám người chúng tôi. Tôi rùng mình, xoay người tìm kiếm.

``A-Ai vậy?''

Khi tôi chớp mắt, chúng hiện ra. Ngay sau lưng Kaoru-san và Yuka-san, những người vẫn đang trò chuyện một cách bình thường.

Inari. Touka.

\img{e2-1i}

Bóng ma của hai người, ánh mắt sắc nhọn như hai viên đạn, liếc thẳng về phía tôi. Giọng nói ấy - chậm, thấp, chẳng để tôi phản bác - chính là của họ.

Khóe môi họ giật nhẹ, như đang cảnh báo. Như nhắc tôi nhớ về cái ngày ám ảnh bắt nguồn cho mọi sự kiện kia.

Tôi không thể cất lời. Chỉ biết chết lặng, trong khi hai bóng ma dần dần trôi ra xa, mờ đi rồi biến mất hoàn toàn.

Rốt cuộc thì\dots\ tại sao họ lại ở đây? Ở thời kỳ này? Đáng lẽ, họ là những người mà chúng tôi sẽ gặp trong tương lai chứ?

Đầu tôi quay cuồng, không thể hiểu nổi những gi đang diễn ra trước mắt. Hàng đống câu hỏi cứ thế mọc lên, nhưng chẳng có cái nào trong số chúng có lời giải đáp.

``\dots-san? Saikyou-san? Cậu ổn chứ?''

Giọng của Yuka-san vang lên, khe khẽ mà như gấp gáp. Ngay sau đó là Akane-san, đôi mắt lo lắng: ``Saikyou-chan\dots\ mặt cậu tái mét hết rồi kìa.''

Tôi giật mình, thoát khỏi khoảng trống vô hồn khi bóng dáng Inari và Touka tan biến. Không chỉ hai người đó, mà ánh mắt của tất cả quanh tôi đều đang đổ dồn về phía tôi: Kaoru-san cau mày, Saya-san và Miku-san nhìn tôi với vẻ thăm dò, Yuka-san thì lặng lẽ nhíu mày, còn Nii-chan cũng không giấu nổi sự ái ngại. Anh ấy thừa hiểu tôi vừa trải qua điều gì, nhưng ở đây thì không thể nói ra.

Cô bạn tóc ngắn nhấp hộp sữa, giọng chắc nịch: ``Quả thật là cậu không ổn rồi\dots\ Để tôi đưa cậu về nhé?''

Tôi vội lắc đầu, xua tay, cố tỏ ra ổn: ``Không sao đâu\dots\ Tôi vẫn chịu được. Mọi người không phải lo đâu.''

Saya tiến lên một bước, giọng mềm mại nhưng không kém phần cứng rắn: ``Nhưng mà, sức khỏe là quan trọng nhất mà\~{} Nếu cảm thấy mệt, tốt nhất nên nghỉ ngơi thì hơn.''

Miku cũng lên tiếng, tay vẫn cầm thanh sô-cô-la dang dở: ``Ừ. Không nên cố quá đâu. Chúng ta có thể tiếp tục vào dịp khác được mà.''

Tôi định phản bác, nhưng Nii-chan đặt tay lên vai tôi. Ánh mắt anh kiên quyết, như muốn tôi phải nghe họ. Reiko cũng ghì lấy tay tôi, giọng dịu lại: ``Onee-chan, để bọn em dìu chị về nhé.''

Cuối cùng, tôi chẳng còn cớ gì để từ chối. Tôi gật nhẹ, để mặc hai người dẫn mình đi. Trước khi rời, anh tôi cúi đầu với cặp bạn thân, những người mà chúng tôi gặp đầu tiên: ``Ngày mai bọn tôi sẽ đến đúng điểm hẹn. Xin lỗi vì đã làm phiền mọi người.''

``Ừm, không sao đâu! Nhớ giữ sức khỏe nhé!'' Akane-san vẫy tay, nụ cười vẫn rạng rỡ nhưng mắt ánh lên chút lo lắng. Yuka-san chỉ khẽ phẩy tay, như muốn giấu đi vẻ nghiêm trọng vừa rồi.

Chúng tôi rời đi, để lại năm cô gái đứng lại với tâm trạng khác nhau.

\pov{-}

Nhìn ba người họ rời đi, cả năm cô gái đều không giấu nổi sự lo lắng trước những người bạn mới làm quen.

Kaoru thở dài, ánh mắt còn vương nỗi bận tâm: ``\dots Cô ấy trông yếu thật đấy. Tôi thấy khó yên tâm quá.''

Saya gật đầu, giọng ngân dài như thở: ``Ừ nhỉ\dots\ Nếu mai mà đi tham quan, mong là Saikyou-san sẽ khỏe hơn.''

Akane ôm cánh tay, lo lắng thật sự:
``Mình thấy cậu ấy cố gắng gượng quá. Không biết có chuyện gì xảy ra với cậu ấy nữa\dots''

Yuka chép miệng, nhưng giọng lại sắc hơn bình thường: ``Chẳng lẽ chỉ là mệt thôi sao? Nhìn kiểu gì cũng không giống.''

Miku im lặng một thoáng, rồi lên tiếng, điềm đạm: ``Dù thế nào thì, mai mình cũng sẽ để ý đến ba người đó. Nhất là Saikyou-san.''

Cả nhóm gật đầu, dù mỗi người một suy nghĩ, chẳng ai thực sự an lòng.

\pov{Haragen Saikyou}

Khi bóng năm người kia đã khuất dần sau hàng cây, ba anh em chúng tôi mới đồng loạt thở phào. Không đơn thuần chỉ vì đã giấu được họ, mà còn vì không muốn kéo người ngoài vào trong mớ hỗn loạn này.

Sẽ ra sao nếu ba chúng tôi nói toẹt hết những gì mà chúng tôi đang nghĩ? Tôi không rõ như thế nào, nhưng chắc chắn sẽ chẳng một ai trong số họ tin cả. Vậy nên, giấu giếm họ là tốt nhất.

Reiko thở dài: ``Chúng ta\dots\ lại rơi vào chuyện đó một lần nữa rồi, đúng không?''

Nii-chan gật, mắt nhìn xa xăm:
``Ngôi đền Aogami\dots\ nơi mà chúng ta đã gặp Shino-san\dots\ và cũng là nơi Sora cầu nguyện sau trận sạt lở. Không thể là trùng hợp được.''

Tôi cắn môi, cảm thấy gai ốc nổi khắp người: ``Rồi những lời đồn kia nữa\dots\ họ đều mất tích ở gần ga. Chính nơi mà chúng ta vừa bước ra. Không lẽ mọi thứ đều xoay quanh ga Aogami?''

Em út rùng mình, siết chặt tay áo:
``Em không muốn tin đâu\dots\ nhưng cứ như là vòng lặp kia đang kéo chúng ta trở lại.''

Anh tôi nhìn thẳng vào tôi, giọng thấp: ``Saikyou, em cũng cảm nhận được phải không? Những cái cảm giác ấy\dots\ thân thuộc tới lạ.''

Tôi khẽ gật, môi run run:
``\dots Ừ. Cả hai chúng ta đều đã thấy. Mà mỗi lần nó đến, nó\dots\ nó không đơn thuần chỉ là cảm giác. Nó kéo cả ký ức lẫn nỗi sợ trở lại\dots''

Khi nhớ lại những gì đã xảy ra xoay quanh thị trấn này, cả ba chúng tôi đều không khỏi rùng mình. Những tiếng thì thầm ma quái, những bóng đen lướt qua trong đêm, và cả những ký ức đau đớn cứ quẩn quanh trong tâm trí. Tôi có thể thấy rõ Reiko đang siết chặt tay áo, còn Nii-chan thì cố kìm nén nỗi sợ đang dâng lên.

``Chúng ta phải hết sức cẩn thận,'' tôi lên tiếng, giọng trầm xuống. ``Không biết được điều gì đang chờ đợi phía trước.''

Cả hai người chỉ có thể gật đầu, ánh mắt đầy lo lắng.

Chúng tôi im lặng bước đi, những cơn gió đông tiếp tục thổi thẳng vào ba chúng tôi. Dường như cả thị trấn đang thì thầm về những vòng lặp mà chỉ có chúng tôi mới nhớ được, mới biết được, và mới trải nghiệm trực tiếp được.

\ct

Không ai trong chúng tôi lên tiếng. Ba giờ chiều, nắng vẫn còn gay gắt, nhưng cái lạnh lẽo bám sau gáy vẫn không tan biến. Dù bề ngoài tôi cố tỏ ra bình thường, nhưng trong lòng vẫn bị trói chặt bởi mớ cảm giác khó hiểu: ký ức xa lạ vừa bị đè nén, và cả những ánh mắt vô hình dõi theo. Tôi tự nhủ rằng mình phải bước tiếp, phải giả vờ như chẳng có gì xảy ra.

Cuối cùng, chúng tôi dừng lại trước một căn nhà quen thuộc. Ba ánh mắt cùng hướng về cánh cổng gỗ nâu sẫm, nơi vòm cây che phủ tạo nên bóng mát âm u dưới ánh nắng chiều. Căn nhà này, theo trí nhớ của tôi, chính là ``nhà mình.''

Nii-chan thở lấy một hơi dài, liếc sang tôi. Tôi chỉ đáp lại bằng một cái gật rất khẽ. Bên cạnh, Reiko cũng thở phào nhẹ nhõm khi chúng tôi cuối cùng cũng đi được tới nơi cần đến.

``Chúng ta\dots\ về rồi.''

Nii-chan khẽ thốt ra, như để tự trấn an bản thân, mặc dù tôi biết sâu trong lòng, anh không có cảm giác như vậy.

Bàn tay anh chạm vào tay nắm cửa. Tôi thấy anh khẽ rùng mình, có lẽ vì cái lạnh kim loại truyền ngược lên. Nhưng cánh cửa rồi cũng mở ra, để lộ khoảng không tối mờ bên trong.

Cả ba chúng tôi bước vào.

Âm thanh cánh cửa khép lại vang vọng khắp căn nhà, như một dấu chấm lửng cho chuyến đi kỳ lạ vừa rồi. Nhưng đồng thời, nó cũng giống như lời báo hiệu cho một điều gì đó còn đang chờ chúng tôi phía trước.

Trong màn im lặng bao trùm, tôi không thể ngừng tự hỏi: Liệu đây thật sự là nhà của chúng tôi?

Và\dots\ nếu đúng vậy, tại sao tôi có cảm giác như thể có điều gì đó sai trái khủng khiếp đang ẩn mình ngay sau lớp tường quen thuộc này?

\end{document}